\section{Numerical Results}
\subsection{Code Search}
Now, we describe the code searching process. First, we will pick $L, J, \Gamma$, a group product $\star$, and two groups: one small non-abelian group $H_k$ and one large abelian group $C_m$. The sweep over these hyperparameters can be restricted by the bounds of Section~\ref{sec:bounds}. Furthermore, the group structures and products used should be customized to match physical constraints and logical needs, which we covered in Section~\ref{sec:hardware} and Section~\ref{sec:logic}.

The sampling of the code proceeds in two stages: we first enumerate permutation-inequivalent top ansatzes $\mathcal{F}^H, \mathcal{G}^H$ that satisfy the anti-commutativity pattern given by the active set $\Gamma$. For small $k$, this can be done efficiently and exhaustively. To sample a DPG, we then randomly sample $L$ bottom elements $(c_1,...,c_L)\in C^L_m$, where $m$ is odd. We can estimate the distance of this bottom code quickly, and, whenever $\deg H_k$ is odd as it is for the $S_3$ tops we use, reject samples with $(\deg H_k)\,d_{\mathrm{bottom}}\leq L$, which by Corollary~\ref{cor:quotient-distance} already forces $d\leq L$. To sample a SPG, we need to pay additional care to ensure commutativity. We can either restrict our ansatzes to ones that meet the sufficient condition listed in Lemma~\ref{lemma:spk-commutivity} to allow for the maximum number of AOD compatible elements (it suffices to leave a $2/L$ fraction of the generators non-compatible in the most restricted case); or we can inherit the DPG ansatzes, and restrict the $c\in C_m$ sampled to only the ones satisfying Equation~\eqref{eq:spk-commutivity}. We label the former SPG-\emph{restricted} and the latter SPG-\emph{full}.

Finally, we post-select on girth $\geq 6$ and use \texttt{QDistRnd} to estimate the code distance~\cite{Pryadko_2022}. A \texttt{QDistRnd} run reports the lowest logical weight it happens to hit, and is therefore an \emph{upper} bound on $d$; we consequently spend trials adaptively, escalating only on codes that are still candidates for a large distance. Every sampled code receives an initial pass of $8$ random trials, and is promoted to the next tier --- $1024$, then $10^4$, then $10^5$ trials --- only while the smallest logical weight found so far is at least $6$, then $8$, then $10$. Codes that survive the last tier and lie on the Pareto frontier receive a final pass of $2\times10^5$ trials. We performed distance estimation in this way on more than $2\times10^5$ randomly sampled codes with girth $\geq 6$ and $n\leq 3000$, across the DPG, SPG-restricted, and SPG-full constructions and predominantly with polynomial lifts, with results reported in Figure~\ref{fig:codes}. Two reference families are plotted in black there for comparison: the co-designed Kasai baselines of~\cite{zhao2026ultrahighratequantumerrorcorrection} and the pair-partition CPM codes of~\cite{okada2026pairpartitionconstructionscpmbasedquantum}.

Because an estimate of this kind can only ever be an upper bound, every instance we report below is additionally \emph{certified}: we determine $d$ exactly by exhaustive enumeration, which is feasible in the regime we target, $n\leq 2500$ and $d\leq 16$. This is the certification standard established for this code family in~\cite{okada2026pairpartitionconstructionscpmbasedquantum}; the enumeration used here is our own implementation. It has two independent parts, and both are carried out separately in the $X$ and in the $Z$ sector. The first is \emph{exclusion}: we rule out every non-trivial logical operator of weight strictly below $d$, i.e.\ a complete search of $\ns{H_Z}\setminus\rs{H_X}$ (for the $X$ sector) and of $\ns{H_X}\setminus\rs{H_Z}$ (for the $Z$ sector) through weight $d-1$. In practice the search only has to reach weight $d-2$: every column weight of $H_X$ and $H_Z$ is odd for the instances reported here, so $\textbf{1}^{1\times n}$ lies in each row space, every logical operator has even weight, and the odd weights can be skipped. The second is a \emph{witness}: we exhibit one explicit operator of weight exactly $d$ in each sector and re-verify directly that it is a non-trivial logical. Together the two parts give $d_X$ and $d_Z$ exactly, hence $d=\min(d_X,d_Z)$. Explicit group-ring generators for the certified frontier instances are tabulated in Tables~\ref{tab:frontier-nondual}--\ref{tab:frontier-other}.
\begin{figure*}
    \centering
        \includegraphics[width=.46\linewidth]{Assets/small_codes.pdf}
        \includegraphics[width=0.46\linewidth]{Assets/dual_codes.pdf}
        \caption{Codes discovered in the preliminary search. Both panels plot certified distance $d$ horizontally against code length $n$ vertically, so towards the lower right --- more distance for fewer qubits --- is better.
        (Left): the practical regime, $n\leq2500$, rate $\geq1/2$, $w\leq12$ and girth $\geq6$; one representative, the best instance, is shown for each (girth, $L$, construction-type) class, and the annotated instances are certified frontier codes, listed with their row weight $w$. Marker shape and color distinguish the construction type (DPG, SPG-restricted, SPG-full) and the ZX-duality sector involution of Eq.~\eqref{eq:zx-sector-maps}, if any.
        (Right): ZX-dual GALA codes. Marker shape and color record the duality class (which sector involution of Eq.~\eqref{eq:zx-sector-maps} the instance realizes, $i\mapsto i$, $i\mapsto i+L/4$, or the reflection $i\mapsto -i$) and  girth. Consistently with Proposition~\ref{prop:zx-girth4}, the reflection class is girth $4$. The highlighted subset is the Pareto frontier of those instances that additionally satisfy $w=12$, girth $\geq6$ and rate $\geq1/2$; it contains the self-dual $[[1056,532,12]]$.}
    \label{fig:codes}
\end{figure*}

In addition to rediscovering codes matching known Kasai parameters, we also find:

\begin{enumerate}
    \item barrier-breaking instances with $d>w=12$: the $[[1752, 880, 14]]$, at $76\%$ of the length of the previously known $[[2304, 1156, \leq 14]]$~\cite{zhao2026ultrahighratequantumerrorcorrection}, and the $[[2232, 1120, 16]]$, which is shorter still than that same $n=2304$ baseline while reaching a higher and exactly certified distance;
    \item a $[[720, 360, 12]]$ at $62.5\%$ of the length of the previously known $[[1152, 580, \leq 12]]$~\cite{zhao2026ultrahighratequantumerrorcorrection}; here $d = w = 12$, so this instance does not itself break the weight barrier, but it attains the baseline distance on $5/8$ of the qubits;
    \item short codes reaching $d = 10$: the direct-product $[[480,240,10]]$, the shortest certified $d=10$ instance we found, and the semi-direct-product $[[528,264,10]]$, the shortest such instance among the SPG families;
    \item self-dual codes $[[132,30,12]]$ and $[[1056, 532,12]]$ with fold-transversal Clifford gates. The two sit in different regimes: $[[1056,532,12]]$ has girth $6$ and rate $532/1056 > 1/2$, whereas $[[132,30,12]]$ has girth $4$ and rate $30/132 \approx 0.227$, so the latter is not covered by the girth-$\geq6$, rate-$\geq1/2$ claims made for the other instances and is reported purely for its compactness and its transversal gates.
\end{enumerate}

One further direct-product instance, on $n=2160$ qubits with $k=1084$, has a completed $X$-sector enumeration establishing $d_X=16$ exactly; its $Z$-sector enumeration is not yet complete, so we do not quote it as a certified $[[n,k,d]]$ instance above. Every instance that we do quote carries an exactly certified distance in the sense of the previous paragraph --- all logical operators below $d$ excluded, plus a verified weight-$d$ witness on each side --- so no ``$\leq$'' appears on any GALA parameters that we report as results. The only upper bounds that appear anywhere in this work are the \texttt{QDistRnd} estimates of the uncertified instances in the right panel of Figure~\ref{fig:codes} and the distances quoted for the baselines.


\subsection{Circuit level logical error rates}

We estimate the logical error rate of our two smallest $d=12$ codes and compare them against the Gross code (BB[[144,12,12]])\cite{Bravyi_2024}, the smallest Pair Partition code (PP[[232, 62,12]]) \cite{okada2026pairpartitionconstructionscpmbasedquantum} and the Surface Code at distance 11 and 13\cite{Fowler_2012}. We use a standard error model where measurement, readout, and all gate errors are parametrized by a single phyiscal error $p$. We do not account for idle errors as consistent with prior work \cite{zhao2026ultrahighratequantumerrorcorrection}; indeed, with coherence time on the scale of 10s of seconds, the idle error during gate layer are many orders of magnitude smaller.



We report the probability that any logical qubit fails in cycle, or  $p_L = 1 - (1-p_{fail})^{1/kN_c}$, using the same decoder family on all baslines to have an apples-to-apples comparison. Sampling with Stim \cite{Gidney_2021} and decoding with a two-stage Relay-BP cascade\cite{zhao2026ultrahighratequantumerrorcorrection} on a greedy coloration syndrome extraction circuit, we accumulate $2.9 \times 10^9$ shots and $1.0 \times 10^4 $ logical failures across [[132,30,12]] and [672,336,12]] and four baselines, with the surface code decoded by minimum-weight matching. The result is reported in Figure~\ref{fig:ler}. We perform two fits to the LER obtained: one from the threshold theorem $C(p/p_{th})^{d_{eff}}$ \cite{Fowler_2012} and the other from the conventional polynomial-exponent form $p^{d/2} \exp(c_0 + c_1 p + c_2 p^2)$ \cite{Bravyi_2024}. We find the pseudo-thresholds of the GALA codes is at around $0.38–0.41\%$.
\begin{figure}[h]
    \centering
    \includegraphics[width=0.95\linewidth]{Assets/ler_circuit_fits.pdf}
    \caption{Circuit Level Logical Error Simulations. \emph{Left:} Each solid data point point correspond to a decoder experiment, where the light data points have less than 10 recorded failures and are not used in the fit. The two fits correspond to threshold theorem $C(p/p_{th})^{d_{eff}}$ \cite{Fowler_2012} (solid) and polynomial-exponent form $p^{d/2} \exp(c_0 + c_1 p + c_2 p^2)$ \cite{Bravyi_2024} (dashed). Note the [[132,30,12]] and [[676, 336,12]] GALA codes reache $p_L \sim 5\times 10^{-8}$ at $p_{phys} = 1.5\times 10^{-3}$ and $2\times 10^{-3}$, which is the lowest recorded experiments. \emph{Right:} Extrapolated logical erorr rates at $10^{-3}$. GALA achieves similar LER performances as the BB codes at much lower qubit overhead. }
    \label{fig:ler}
\end{figure}


Extrapolating to $p = 10^{-3}$, [[672,336,12]] reaches on average $5 \times 10^{-11}$ per logical qubit per cycle at an overhead of 3.0 physical qubits per logical qubit, against $2 \times 10^-12$ at an overhead of 24 for the gross code and $4 \times 10^{-12}$ at an overhead of 5.6 for the pair-partition code. The rate-1/2 block therefore gives up roughly 1.4 decades of logical error rate in exchange for an eightfold reduction in qubit overhead relative to gross. 

We note that $d_{eff}$ is estimated to be $9.5$ and $13.5$ for the [[132,30,12]] and [672,336,12]] codes in this range, indicating near-threshold waterfall behavior granting us additional error correction efficient, as similarly observed in prior work \cite{kasai2026breakingorthogonalitybarrierquantum,zhao2026ultrahighratequantumerrorcorrection}. For this reason, our extrapolation should be interpreted as a rough estimate, most accurate for GALA codes and surface codes since we have data points at $p = 1.5^10^3$ and $p = 2\times 10^{3}$ near $10^{-3}$. Our experiment on the [[132,30,12]] code at $p = 10^{-3}$ using the full cascade including MILP \cite{zhao2026ultrahighratequantumerrorcorrection} timed out at $5\times 10^6$ shots \emph{without a single error}, suggesting an estimated ceiling of ~$10^{-8}$, which is consistent with our extrapolation.  








